\chapter{Hepatectomy}

\section*{Introduction \& Clinical Indications}
Hepatectomy is a major surgical procedure involving the partial or complete removal of the liver, or more commonly, a specific segment or lobe of the liver. This intervention is primarily undertaken to excise malignant or benign tumors, treat severe traumatic injuries, or manage other localized liver pathologies that are unresponsive to non-surgical therapies. The anatomical site of the surgery is the liver, a vital organ located in the upper right quadrant of the abdomen, crucial for metabolism, detoxification, and protein synthesis. Clinically, hepatectomy is a cornerstone treatment for primary liver cancers like hepatocellular carcinoma and intrahepatic cholangiocarcinoma, as well as for metastatic disease, particularly from colorectal cancer. Its clinical significance lies in offering a potentially curative option for diseases that would otherwise be fatal, demanding meticulous surgical precision and a profound understanding of hepatic anatomy and physiology to preserve adequate functional liver remnant.

\begin{itemize}
  \item Liver resection
  \item Partial hepatectomy
  \item Segmentectomy
  \item Lobectomy of the liver
  \item Wedge resection
\end{itemize}

The fundamental surgical principle of hepatectomy is based on the Couinaud classification, which divides the liver into eight functionally independent segments, each with its own portal triad (hepatic artery, portal vein, bile duct) and hepatic venous drainage. This anatomical understanding allows for precise, oncologically sound resections that minimize blood loss and preserve the viability of the remaining liver. The operative concept involves the meticulous excision of diseased liver parenchyma, typically along these segmental boundaries, while ensuring the future liver remnant (FLR) retains adequate vascular inflow and outflow, as well as biliary drainage. The primary technical goal is to achieve a complete (R0) resection of the pathology, especially for malignancy, with clear surgical margins.

The rationale for hepatectomy is rooted in the liver's remarkable regenerative capacity and its segmental vascular architecture. By selectively ligating or dividing the segmental pedicles (inflow) and hepatic veins (outflow) supplying the diseased portion, surgeons can isolate and remove the target tissue with minimal impact on the remaining healthy liver. This approach minimizes blood loss and bile leakage during parenchymal transection. Preoperative assessment of the FLR volume and function is critical to prevent post-hepatectomy liver failure, sometimes necessitating portal vein embolization or staged resections to induce hypertrophy of the FLR prior to the definitive surgery, thereby optimizing patient safety and oncologic outcomes.

The history of hepatectomy dates back to the late 19th century, with the first successful elective liver resection performed by Carl Langenbuch in 1888 for a benign tumor. Early attempts were fraught with high mortality rates, primarily due to uncontrolled hemorrhage and sepsis. Significant advancements in understanding hepatic anatomy and surgical technique occurred throughout the 20th century, making major resections progressively safer. A pivotal moment came in the 1950s with Claude Couinaud's detailed anatomical descriptions of the liver's segmental vascular and biliary architecture, which revolutionized liver surgery by enabling precise, anatomically guided resections rather than crude wedge excisions.

Further evolution included the widespread adoption of the Pringle maneuver (temporary occlusion of the portal triad) to reduce intraoperative blood loss, and the development of sophisticated parenchymal transection devices such as ultrasonic aspirators (CUSA) and water-jet dissectors. The 1980s and 1990s witnessed a dramatic increase in the safety and scope of hepatectomy, driven by improved critical care, advanced anesthesia techniques, and a deeper understanding of liver regeneration. More recently, the advent of minimally invasive approaches, including laparoscopic and robotic hepatectomy, has further refined the procedure, offering reduced postoperative pain, shorter hospital stays, and faster recovery for selected patients, marking a continuous evolution towards safer and more effective liver surgery.

\begin{itemize}
  \item To remove primary malignant liver tumors such as hepatocellular carcinoma (HCC) in patients with compensated liver function or non-cirrhotic livers, often guided by Milan criteria for transplant eligibility or Barcelona Clinic Liver Cancer (BCLC) staging.
  \item To resect intrahepatic cholangiocarcinoma, a malignant tumor arising from the bile ducts within the liver, where complete surgical excision offers the only chance for cure.
  \item To treat colorectal liver metastases (CRLM), which are common and potentially curable with resection in selected patients, often after multidisciplinary team discussion and neoadjuvant chemotherapy.
  \item To remove metastases from other primary cancers, including neuroendocrine tumors, breast cancer, or sarcomas, when the disease is confined to the liver and amenable to complete resection.
  \item To excise large, symptomatic benign liver tumors such as giant hemangiomas, focal nodular hyperplasia, or adenomas with a high risk of hemorrhage, rupture, or malignant transformation.
  \item To manage severe traumatic liver injuries involving devitalized tissue or ongoing hemorrhage unresponsive to non-operative measures, such as angioembolization.
  \item To remove liver abscesses or parasitic cysts (e.g., echinococcal cysts) that are refractory to medical treatment, percutaneous drainage, or pose a significant risk of rupture.
  \item As part of living donor liver transplantation, where a portion of a healthy donor's liver is resected for transplantation into a recipient, requiring meticulous assessment of donor safety.
\end{itemize}

Hepatectomy is often considered a primary, potentially curative treatment for resectable liver malignancy, particularly for early-stage hepatocellular carcinoma or isolated colorectal liver metastases. In these scenarios, it represents the first-line surgical intervention aimed at complete tumor eradication, offering the best chance for long-term survival. However, it can also serve as a secondary or salvage procedure within a comprehensive oncologic strategy. For instance, in patients with metastatic disease, hepatectomy may be performed after neoadjuvant chemotherapy has reduced tumor burden or size, making previously unresectable lesions amenable to surgery. It can also be a salvage option following failed local ablative therapies, such as radiofrequency ablation or transarterial chemoembolization, if the disease remains localized and resectable. Furthermore, hepatectomy frequently functions as an adjunct procedure within a multi-modality treatment plan, combined with chemotherapy, radiation, or targeted therapies, to optimize oncologic outcomes and improve patient prognosis.

\section*{Relevant Surgical Anatomy}
The liver is the largest solid organ in the human body, situated in the right upper quadrant of the abdomen, beneath the diaphragm, and protected by the lower ribs. It is broadly divided into a larger right lobe and a smaller left lobe by the falciform ligament anteriorly, which extends to the diaphragm. Posteriorly, the liver is attached to the diaphragm by the coronary and triangular ligaments, leaving a "bare area" that is not covered by peritoneum. The porta hepatis, located on the inferior surface of the liver, serves as the entry point for the hepatic artery and portal vein, and the exit point for the common hepatic duct. The inferior vena cava (IVC) runs in a groove on the posterior surface of the liver, receiving the hepatic veins.

The functional anatomy of the liver is best understood through the Couinaud classification, which divides the liver into eight functionally independent segments. Each segment has its own vascular inflow (branches of the portal vein and hepatic artery) and biliary drainage, allowing for precise, anatomically guided resections. The portal vein and hepatic artery branches, along with the bile ducts, form the portal triads that run centrally within each segment. The hepatic veins (right, middle, and left) run intersegmentally, draining blood from the segments into the IVC. The right hepatic vein typically drains segments V, VI, VII, and VIII. The middle hepatic vein drains segments IV, V, and VIII, while the left hepatic vein drains segments II, III, and part of IV. Segment I (caudate lobe) has independent vascular supply and drainage directly into the IVC.

Lymphatic drainage of the liver primarily follows the portal triads to lymph nodes in the porta hepatis, which then drain to celiac and peripancreatic nodes. The nerve supply to the liver is derived from the vagus nerve (parasympathetic) and the celiac plexus (sympathetic). Key surgical landmarks include the porta hepatis, the gallbladder fossa, the IVC, and the major hepatic veins, which must be carefully identified and preserved or controlled during resection. Understanding these intricate anatomical relationships is paramount for safe and effective hepatectomy, minimizing blood loss and preserving the function of the remaining liver.

\section*{Preoperative Workup \& Preparation}
Thorough preoperative preparation is crucial for optimizing patient outcomes following hepatectomy, minimizing risks, and ensuring adequate future liver remnant function. This typically involves a comprehensive assessment of the patient's overall health, liver function, and the extent of the disease, often guided by a multidisciplinary team. Key components include detailed medical and cardiac evaluation to ensure fitness for major surgery, extensive imaging to precisely map the tumor and assess liver volume, and optimization of nutritional status. Patients are also counselled on the procedure, risks, and expected recovery.

Specific preparatory steps include:
\begin{itemize}
  \item \textbf{Laboratory Tests:} Complete blood count (CBC), liver function tests (LFTs), coagulation panel (INR, PTT), renal function tests (creatinine, BUN), electrolytes, and tumor markers (e.g., AFP for HCC, CEA for CRLM).
  \item \textbf{Imaging Studies:} Triphasic computed tomography (CT) or magnetic resonance imaging (MRI) of the abdomen with contrast is essential to define the tumor's size, location, relationship to major vascular structures, and to assess the volume and quality of the future liver remnant (FLR). A chest CT and often a PET-CT scan are performed for comprehensive staging in malignancy.
  \item \textbf{Cardiac and Anesthesia Clearance:} Electrocardiogram (ECG), echocardiogram, and pulmonary function tests (PFTs) may be required, especially for older patients or those with comorbidities, to ensure fitness for general anesthesia and major surgery.
  \item \textbf{Medication Adjustments:} Anticoagulants and antiplatelet medications are typically discontinued several days prior to surgery, under medical supervision, to minimize bleeding risk.
  \item \textbf{Prophylactic Antibiotics:} Broad-spectrum intravenous antibiotics are administered shortly before the incision to reduce the risk of surgical site infection.
  \item \textbf{Informed Consent:} A detailed discussion with the patient regarding the risks, benefits, alternatives, and potential complications of hepatectomy is conducted, and informed consent is obtained.
  \item \textbf{Future Liver Remnant Augmentation:} For patients requiring extensive resections, portal vein embolization (PVE) or associating liver partition and portal vein ligation for staged hepatectomy (ALPPS) may be performed weeks prior to surgery to induce hypertrophy of the FLR, thereby reducing the risk of post-hepatectomy liver failure.
\end{itemize}

\textbf{Absolute Contraindications:}
\begin{itemize}
  \item Unresectable disease, characterized by diffuse bilobar involvement, extensive invasion of major vascular structures (e.g., main portal vein, inferior vena cava), or widespread extrahepatic metastases.
  \item Severe liver dysfunction, typically Child-Pugh Class C cirrhosis, which indicates a high risk of postoperative liver failure.
  \item Inadequate future liver remnant (FLR) volume and function, even after attempts at FLR augmentation (e.g., portal vein embolization).
  \item Severe uncorrectable coagulopathy, which significantly increases the risk of life-threatening hemorrhage.
  \item Severe cardiopulmonary disease or other uncorrectable comorbidities that render the patient unfit for major surgery and general anesthesia.
\end{itemize}

\textbf{Relative Contraindications and Precautions:}
\begin{itemize}
  \item Child-Pugh Class B cirrhosis, where the risk of liver failure is elevated but may be acceptable in highly selected cases with small resections.
  \item Significant portal hypertension with esophageal varices or ascites, which increases the risk of bleeding and postoperative complications.
  \item Active systemic infection or sepsis, which should be treated and resolved prior to elective hepatectomy.
  \item Prior extensive abdominal surgery, which may complicate liver mobilization and increase operative time due to adhesions.
  \item Precaution: Meticulous hemostasis throughout the procedure is paramount to prevent significant blood loss.
  \item Precaution: Careful identification and preservation of bile ducts and major vascular structures are essential to avoid iatrogenic injury.
  \item Precaution: Intraoperative ultrasound is critical for precise tumor localization and defining safe resection margins, especially in complex cases.
\end{itemize}

\section*{Surgical Technique \& Procedure}
Hepatectomy is performed under general anesthesia, often supplemented with epidural analgesia for postoperative pain control. The patient is typically positioned supine, sometimes with the right arm tucked, to allow for optimal surgical access. The choice between an open, laparoscopic, or robotic approach depends on the tumor's size, location, complexity, the patient's liver function, and the surgeon's expertise. The procedure generally follows these key steps:

\begin{enumerate}
  \item \textbf{Incision and Exploration:} For open surgery, a right subcostal (Kocher) incision, a bilateral subcostal (Mercedes) incision, or a midline incision may be used. Laparoscopic or robotic approaches involve multiple small port incisions. The abdomen is thoroughly explored to rule out any unexpected extrahepatic disease or peritoneal carcinomatosis.
  \item \textbf{Liver Mobilization:} The liver is extensively mobilized by dividing its ligamentous attachments (falciform, triangular, coronary ligaments) to allow for adequate exposure, manipulation, and access to the retrohepatic inferior vena cava.
  \item \textbf{Vascular and Biliary Control:} The hepatic inflow (hepatic artery and portal vein branches) and outflow (hepatic veins) to the segment or lobe targeted for resection are identified and controlled. This may involve selective clamping of individual pedicles or the Pringle maneuver (temporary clamping of the entire portal triad) to minimize blood loss during parenchymal transection.
  \item \textbf{Resection Plane Identification:} Intraoperative ultrasound is routinely used to precisely localize the tumor, define the resection margins, and identify key vascular and biliary structures within the liver parenchyma. The planned transection line is marked on the liver capsule using electrocautery.
  \item \textbf{Parenchymal Transection:} The liver parenchyma is carefully divided along the predetermined resection plane. Various techniques are utilized, including the crush-clamp method, ultrasonic aspirators (CUSA), water-jet dissectors, radiofrequency devices, or surgical staplers. As the parenchyma is divided, small vessels and bile ducts are meticulously ligated, clipped, or coagulated to ensure hemostasis and prevent bile leakage.
  \item \textbf{Specimen Removal and Hemostasis:} Once the segment or lobe is completely separated, the specimen is removed. The raw liver surface is meticulously inspected for any bleeding points or bile leaks, which are then controlled with sutures, clips, or electrocautery. The integrity of the remaining vascular and biliary structures is confirmed.
  \item \textbf{Drain Placement and Closure:} A surgical drain, typically a closed-suction drain, may be placed near the resection site to monitor for bile leak or fluid collection. The abdominal wall is then closed in layers, and the skin is closed with sutures or staples.
\end{enumerate}

\section*{Complications \& Risks}
Hepatectomy is a major surgical procedure associated with a range of potential risks, which can be broadly categorized into general surgical risks and procedure-specific complications.

\textbf{General Surgical Risks:}
\begin{itemize}
  \item \textbf{Bleeding:} Significant intraoperative hemorrhage requiring blood transfusions, or postoperative hematoma formation.
  \item \textbf{Infection:} Surgical site infection (wound infection), intra-abdominal abscess, pneumonia, or sepsis.
  \item \textbf{Anesthesia Complications:} Adverse reactions to anesthetic agents, cardiac events (myocardial infarction, arrhythmias), pulmonary complications (atelectasis, acute respiratory distress syndrome), or stroke.
  \item \textbf{Thromboembolism:} Deep vein thrombosis (DVT) in the legs, which can lead to pulmonary embolism (PE), a potentially life-threatening condition.
\end{itemize}

\textbf{Procedure-Specific Risks:}
\begin{itemize}
  \item \textbf{Post-hepatectomy Liver Failure (PHLF):} A severe complication where the remaining liver tissue is insufficient to maintain normal function, leading to jaundice, coagulopathy, and encephalopathy. This is a major cause of mortality.
  \item \textbf{Bile Leak/Biliary Fistula:} Leakage of bile from the raw liver surface or a ligated bile duct, which can lead to intra-abdominal fluid collections, abscesses, or peritonitis.
  \item \textbf{Pleural Effusion/Ascites:} Accumulation of fluid in the chest cavity or abdominal cavity, respectively, often requiring drainage.
  \item \textbf{Post-hepatectomy Hemorrhage:} Bleeding from the raw liver surface or ligated vessels in the postoperative period, potentially requiring re-operation.
  \item \textbf{Renal Insufficiency:} Kidney dysfunction, which can be exacerbated by blood loss, hypotension, or certain medications.
  \item \textbf{Pancreatitis:} Inflammation of the pancreas, a rare but serious complication, possibly due to surgical manipulation or ischemia.
  \item \textbf{Incisional Hernia:} Weakness in the abdominal wall at the site of the incision, leading to protrusion of abdominal contents.
  \item \textbf{Recurrence of Primary Disease:} For malignant conditions, there is always a risk that the cancer may return in the liver or elsewhere in the body, even after a successful resection.
  \item \textbf{Mortality:} Despite advancements, hepatectomy carries a small but significant risk of death, particularly in patients with underlying liver disease or extensive resections.
\end{itemize}

\section*{Postoperative Care \& Recovery}
Immediately after hepatectomy, patients are typically transferred to the Post-Anesthesia Care Unit (PACU) for close monitoring of vital signs, pain levels, and fluid balance. Depending on the extent of surgery and patient comorbidities, they may then be admitted to an intensive care unit (ICU) or a high-dependency unit for the first 24-48 hours. During this critical period, pain management is a priority, often utilizing epidural catheters or patient-controlled analgesia (PCA). Close attention is paid to liver function tests (bilirubin, INR, albumin), renal function, and electrolyte balance, with early detection of any signs of liver dysfunction. Early mobilization, deep breathing exercises, and incentive spirometry are actively encouraged to prevent pulmonary complications. Drain output is meticulously monitored for volume and character, particularly for signs of bile leak or hemorrhage.

Over the first 1-2 weeks, patients gradually progress through their recovery timeline. The typical hospital stay ranges from 5 to 10 days, depending on the complexity of the surgery and the individual patient's recovery trajectory. Drains are usually removed when output is minimal and non-bilious. The diet is advanced from clear liquids to solids as tolerated, and pain management transitions to oral medications. Patients are encouraged to increase their activity levels progressively, aiming for independence in daily activities. Long-term recovery involves a period of fatigue that can last several weeks to months. The liver remnant undergoes significant regeneration, typically reaching its original volume within 6-12 months. Patients are advised to avoid heavy lifting or strenuous abdominal activities for 6-8 weeks to allow for complete abdominal wall healing and to prevent incisional hernia.

During hepatectomy, the surgeon meticulously documents several key findings. These include the precise location, size, and gross appearance of the lesion(s), their relationship to major vascular and biliary structures, and the estimated volume of the resected liver segment or lobe. The presence of any satellite lesions, capsular invasion, or adherence to adjacent organs is also noted. Intraoperative ultrasound findings, which guide the resection plane and confirm tumor clearance, are crucial surgical findings. The completeness of the gross resection, aiming for clear margins, is a primary intraoperative goal.

The resected liver specimen is then sent for detailed histopathological examination, which forms the cornerstone of the definitive diagnosis and guides subsequent management. The pathology report provides critical information: confirming the tumor type (e.g., hepatocellular carcinoma, cholangiocarcinoma, metastatic adenocarcinoma), its histological grade, and precise dimensions. It also assesses for microvascular invasion, perineural invasion, and the presence of satellite nodules. Most importantly, the pathologist determines the margin status: an R0 resection signifies no tumor cells at the surgical margin, indicating a complete microscopic removal; an R1 resection denotes microscopic tumor cells at the margin; and an R2 resection indicates macroscopic residual tumor. The presence of positive margins (R1 or R2) has significant prognostic implications and may necessitate further adjuvant therapy or re-resection.

A normal and successful postoperative course following hepatectomy is characterized by a predictable progression towards recovery and the absence of major complications. Patients typically experience a gradual resolution of incisional pain, which is well-managed with a combination of intravenous and oral analgesics, transitioning to oral medications within a few days. Liver function tests, which may initially be elevated due to surgical stress, are expected to normalize progressively over days to weeks, indicating successful regeneration and function of the remaining liver remnant. Patients should be able to tolerate a regular diet within 3-5 days post-surgery, with bowel function returning to normal. Early ambulation is encouraged, leading to a steady increase in activity levels. Drains, if placed, will show decreasing, non-bilious, and non-hemorrhagic output, allowing for their removal typically within the first week. The surgical wound should heal without signs of infection, and the patient's overall energy levels will gradually improve, allowing a return to light activities within 2-4 weeks and full recovery, including strenuous activity, by 2-3 months.

Postoperative care and long-term follow-up are essential after hepatectomy to monitor for recovery, detect potential complications, and surveil for disease recurrence, especially in oncologic cases. This comprehensive approach ensures optimal patient outcomes and timely intervention if issues arise.

Key aspects of postoperative care and follow-up include:
\begin{itemize}
  \item \textbf{Postoperative Clinic Visits:} Initial visits are typically scheduled 1-2 weeks after discharge for wound check, drain removal (if still present), and general recovery assessment. Subsequent visits at 4-6 weeks and then periodically evaluate overall healing and functional status.
  \item \textbf{Oncologic Surveillance:} For malignant conditions, regular follow-up is essential. This usually involves clinic visits every 3-6 months for the first 2-3 years, then annually, depending on the specific cancer type and risk of recurrence.
  \item \textbf{Imaging Studies:} Cross-sectional imaging, such as CT or MRI of the abdomen and pelvis, is performed periodically (e.g., every 3-6 months) to monitor for local recurrence or new metastases. A chest CT may also be included for systemic disease surveillance.
  \item \textbf{Laboratory Tests:} Serial blood tests, including liver function tests and tumor markers (e.g., CEA for colorectal cancer, AFP for HCC), are conducted to detect early signs of recurrence or liver dysfunction.
  \item \textbf{Activity Resumption Guidance:} Patients receive specific instructions on gradually increasing physical activity, with advice to avoid heavy lifting or strenuous abdominal exercises for several weeks to months to prevent incisional hernia.
  \item \textbf{Warning Signs Education:} Patients are educated on warning signs that warrant immediate medical attention, such as persistent fever, chills, severe abdominal pain, jaundice, dark urine, light-colored stools, or significant abdominal distension, which could indicate complications like infection or bile leak.
\end{itemize}

\section*{Outcomes \& Prognosis}
The epidemiology of hepatectomy is intrinsically linked to the incidence and prevalence of the underlying liver diseases it treats, primarily liver cancers. Hepatocellular carcinoma (HCC), the most common primary liver cancer, has a global incidence that varies significantly, with higher rates in regions endemic for chronic hepatitis B and C infections, as well as in populations with high rates of alcohol-related liver disease and non-alcoholic fatty liver disease (NAFLD). HCC typically affects older adults, with a male predominance. Intrahepatic cholangiocarcinoma, another primary liver malignancy, is less common but its incidence is rising in some Western countries. Colorectal liver metastases (CRLM) are highly prevalent, affecting 25-50\% of patients with colorectal cancer during their disease course, making them the most common indication for liver resection in many Western centers. The age distribution for CRLM reflects that of colorectal cancer, typically affecting individuals over 50. The overall morbidity rate following hepatectomy ranges from 20-40\%, encompassing complications such as bile leak, infection, and liver failure. Mortality rates have significantly decreased over the past decades, now typically ranging from 1-5\% in high-volume specialized centers, reflecting advancements in surgical technique, anesthesia, and perioperative care.

The outcomes and prognosis following hepatectomy are highly variable, depending on the underlying pathology, the extent of liver resection, the patient's overall health, and the presence of underlying liver disease. For primary liver cancers like hepatocellular carcinoma (HCC), 5-year survival rates after curative resection range from 50-70\% for early-stage disease, but recurrence remains a significant challenge, with rates as high as 70\% within 5 years. For colorectal liver metastases (CRLM), 5-year survival rates after resection typically range from 40-60\%, with prognostic factors including the number and size of metastases, margin status, and response to neoadjuvant chemotherapy. Landmark studies, such as those evaluating the role of perioperative chemotherapy in CRLM (e.g., EORTC 40983 trial), have significantly shaped current treatment algorithms. Recurrence rates for CRLM are also substantial, often exceeding 50\%. Key prognostic factors across all indications include the tumor biology (grade, vascular invasion), the completeness of resection (R0 margin status), the quality and volume of the future liver remnant, and the presence of any underlying liver disease (e.g., cirrhosis). Most patients who recover from hepatectomy can expect to return to a good quality of life, although long-term surveillance is crucial to monitor for recurrence and manage any late complications. Functional outcomes are generally good, with the regenerating liver typically able to meet metabolic demands.

\section*{References}
\begin{enumerate}
  \item MedlinePlus. Liver resection. U.S. National Library of Medicine; 2025.
  \item Schwartz's Principles of Surgery. 11th ed. McGraw-Hill Education; 2019.
  \item American Association for the Study of Liver Diseases (AASLD). Practice Guidelines for the Management of Hepatocellular Carcinoma; 2023.
  \item National Comprehensive Cancer Network (NCCN). Clinical Practice Guidelines in Oncology: Colon Cancer; 2024.
\end{enumerate}

\clearpage
